\section{总结与延伸}

\begin{figure}[H]
\centering
\includegraphics[width=0.86\textwidth]{slide-058.jpg}
\caption{Slide 58：Scaling in the wild recap：实践挑战包括架构超参、optimizer 超参和拟合大 Chinchilla sweep 的 compute。}
\figsource{CS336 Spring 2026 Lecture 11 官方幻灯片。}
\end{figure}

\begin{importantbox}{读图：Slide 58 的最终框架}
真实 scaling 的挑战有三类：模型架构超参如何设，optimizer/LR/batch 如何随 scale 变，拟合 compute-optimal sweep 本身如何不花掉太多 compute。解决方案对应三类：muP 或架构假设降低搜索空间，empirical LR/batch scaling 拟合超参，WSD/IsoFLOP/joint fit 降低 Chinchilla sweep 成本。
\end{importantbox}

\subsection{本章小结}

Lecture 11 把 scaling laws 从干净公式带到真实模型训练报告中。MiniCPM 展示了 muP 和 WSD 如何降低 scaling 成本；DeepSeek 展示了直接 batch/LR 和 IsoFLOP 拟合；近期模型说明 scaling 的对象已经扩展到 sparsity、MoE active parameters、architecture 和 downstream 能力；optimizer scaling 和 muP 深入则提醒我们，初始化、LR、batch、optimizer 都可能随 scale 改变。

\begin{importantbox}{最终 takeaways}
\begin{enumerate}
    \item Scaling in practice 的核心不是拟合一条线，而是降低大模型训练前的决策风险。
    \item MiniCPM 的路线是 muP 稳定超参迁移，WSD 降低从头训练网格成本。
    \item DeepSeek 的路线是不依赖 muP，直接做 batch/LR 和 IsoFLOP scaling。
    \item 近期模型把 scaling 扩展到 MoE sparsity、active parameters、architecture 和 downstream metrics。
    \item Optimizer 比较必须做 scale-aware tuning，否则很容易比较错。
    \item muP 有用但不万能；RMSNorm gains、exotic optimizers、strong weight decay 等都可能破坏其假设。
\end{enumerate}
\end{importantbox}

\begin{knowledgebox}{拓展阅读}
\small
\begin{tabular}{L{0.24\textwidth}L{0.34\textwidth}L{0.32\textwidth}}
\toprule
阅读对象 & 带着什么问题读 & 与本讲的连接 \\
\midrule
MiniCPM technical report & muP、WSD 与 scaling computation 分别节省了哪一类搜索成本？ & 对应 Slides 6--18 的完整 recipe。 \\
DeepSeek scaling analysis & 不依赖 muP 时，LR、batch 与 IsoFLOP 最优点如何直接拟合？ & 对应 Slides 19--24 的经验外推路线。 \\
Qwen、Kimi、Hunyuan、LLaMA 3、MiniMax reports & 不同团队选择了哪些 scaling 对象与验收指标？ & 对应 Slides 25--30，比较 dense、MoE 与 downstream scaling。 \\
StepFun scaling study & 公平比较 optimizer 时，为什么必须让 LR 与 batch 一起随 scale 重调？ & 对应 Slides 31--41 的 optimizer scaling 方法论。 \\
muP 与 CerebrasGPT & activation 与 update 的 $\Theta(1)$ 条件如何导出 layer-wise rules？ & 对应 Slides 44--57 的推导与现代组件反例。 \\
Muon scaling studies & 正交化矩阵更新在什么规模、compute 与 tuning 条件下仍有收益？ & 对应 Slides 42--43，避免把单点胜负当成普遍规律。 \\
\bottomrule
\end{tabular}
\end{knowledgebox}

\begin{importantbox}{最小复现实验：先验证迁移，再相信外推}
选择三个宽度递增、架构其余部分尽量固定的小模型。先在最小模型上搜索 learning rate 与 batch size，再分别用 SP 和 muP 将最优设置迁移到后两个模型；每个规模记录训练早期的 activation RMS、update-to-weight ratio、稳定训练步数与最终 validation loss。最后增加一个未参与拟合的中间宽度作为 hold-out：只有当它同时落在 loss 外推曲线内、训练动力学没有 blow up，且迁移超参优于重新粗搜的基线时，才有证据说明 scaling recipe 真正在该架构上成立。
\end{importantbox}

\end{document}
